\DocumentMetadata{
  pdfversion=2.0,
  pdfstandard=ua-2,
  lang=en-US,
  tagging=on,
  tagging-setup={math/setup=mathml-SE,tabsorder=structure}
}
\documentclass[11pt,letterpaper]{article}
\usepackage[margin=0.68in,includefoot]{geometry}
\usepackage{newcomputermodern}
\usepackage{amsmath,amscd,mathtools}
\usepackage{tikz,adjustbox}
\providecommand{\square}{\mdlgwhtsquare}
\usepackage[hidelinks]{hyperref}
\hypersetup{
  pdftitle={Complex Analysis PhD Exam, September 2006},
  pdfauthor={University of Florida Department of Mathematics},
  pdfsubject={Graduate mathematics examination},
  pdfdisplaydoctitle=true,
  bookmarksopen=true
}
\setcounter{secnumdepth}{0}
\setlength{\parindent}{0pt}
\setlength{\parskip}{0.28em}
\setlength{\emergencystretch}{3em}
\setlength{\leftmargini}{2.15em}
\setlength{\leftmarginii}{2.65em}
\setlength{\leftmarginiii}{3em}
\pagestyle{empty}
\raggedbottom
\allowdisplaybreaks
\NewTaggingSocketPlug{block/list/label}{examproblem}{%
  \tagstructbegin{tag=Lbl}\tagstructend%
  \tagstructbegin{tag=\LBody}%
  \tagstructbegin{tag=\ProblemHeadingTag,title-o={\ProblemHeadingText}}%
  \tagmcbegin{tag=\ProblemHeadingTag,actualtext={\ProblemHeadingText}}%
  #2%
  \tagmcend%
  \tagstructend%
}
\newenvironment{examproblems}[2]{%
  \def\ProblemHeadingTag{#2}%
  \AssignTaggingSocketPlug{block/list/label}{examproblem}%
  \begin{enumerate}%
  \setcounter{enumi}{#1}%
  \renewcommand{\labelenumi}{\textbf{\theenumi.}}%
  \setlength{\topsep}{0.35em}%
  \setlength{\partopsep}{0pt}%
  \setlength{\itemsep}{0.48em}%
  \setlength{\parsep}{0.18em}%
}{\end{enumerate}}
\newenvironment{examsubparts}{%
  \AssignTaggingSocketPlug{block/list/label}{default}%
  \begin{enumerate}%
  \setlength{\topsep}{0.18em}%
  \setlength{\partopsep}{0pt}%
  \setlength{\itemsep}{0.16em}%
  \setlength{\parsep}{0.1em}%
}{\end{enumerate}}
\newenvironment{examinstructions}{%
  \par\begingroup\itshape
}{%
  \par\endgroup\vspace{0.18em}
}
\makeatletter
\renewcommand\subsection{\@startsection{subsection}{2}{0pt}%
  {-1.25ex plus -.25ex minus -.1ex}%
  {0.55ex plus .1ex}%
  {\normalfont\large\bfseries}}
\renewcommand{\@maketitle}{%
  \newpage\null\vskip -1em%
  {\footnotesize\bfseries
    \href{https://www.ufl.edu/}{University of Florida}, \href{https://math.ufl.edu/}{Department of Mathematics}\par}%
  \vskip -0.5em%
  \tagstructbegin{tag=H1,title={Complex Analysis PhD Exam, September 2006}}%
  \tagpdfparaOff%
  \tagmcbegin{tag=H1}%
    {\LARGE\bfseries\href{https://gma.math.ufl.edu/exams/complex-analysis-phd/}{Complex Analysis PhD Exam}, September 2006\par}%
  \tagmcend%
  \tagpdfparaOn%
  \tagstructend%
  \vskip 0.25em%
  \tagmcbegin{artifact}%
    \hrule height 1pt%
  \tagmcend%
  \vskip 1em%
}
\makeatother
\title{Complex Analysis PhD Exam, September 2006}
\author{}
\date{}
\begin{document}
\maketitle
\thispagestyle{empty}
\begin{examinstructions}
Give complete proofs and computations. Partial credit will be given where justified. In the following, \(\mathbb{C}\) denotes the set of complex numbers and \(D=\{z\in\mathbb{C}\mid |z|<1\}\).
\end{examinstructions}
\begin{examproblems}{0}{H2}
\def\ProblemHeadingText{Problem 1.}
\item
Evaluate the integral
\[
\int_0^\infty \frac{\log x}{(1+x^2)^2}\,dx.
\]
\def\ProblemHeadingText{Problem 2.}
\item
This problem has two parts.
\begin{examsubparts}
\item[\textnormal{(a)}]
Let~\(f\) be analytic at \(z\in\mathbb{C}\). Prove that
\[
\left(\frac{\partial^2}{\partial x^2}+\frac{\partial^2}{\partial y^2}\right)|f(z)|^2=4|f'(z)|^2.
\]
\item[\textnormal{(b)}]
Let \(f_1,f_2,\ldots,f_n\) be analytic in the region~\(G\). Prove that \(|f_1|^2+|f_2|^2+\cdots+|f_n|^2\) is harmonic on~\(G\) if and only if \(f_k\) is constant, for all \(k=1,\ldots,n\).
\end{examsubparts}
\def\ProblemHeadingText{Problem 3.}
\item
Let \(\{f_n\}\) be a sequence of functions, each analytic in the open set~\(G\), which converges to~\(f\) uniformly on all compact subsets of~\(G\). Prove that~\(f\) is analytic in~\(G\).
\def\ProblemHeadingText{Problem 4.}
\item
Let \(P(z)=z^7+z^5+5z^3+1\). Find the number of zeros of~\(P\), counted according to their multiplicities, in
\begin{examsubparts}
\item[\textnormal{(a)}]
\(\{z\in\mathbb{C}\mid |z|<1\}\)
\item[\textnormal{(b)}]
\(\{z\in\mathbb{C}\mid 1<|z|<2\}\)
\item[\textnormal{(c)}]
\(\{z\in\mathbb{C}\mid |z|>0\}\)
\end{examsubparts}
\def\ProblemHeadingText{Problem 5.}
\item
Let~\(G\) be a region and~\(f\) be analytic in~\(G\). Prove that if \(f(G)\) is a subset of a circle, then~\(f\) is a constant function.
\def\ProblemHeadingText{Problem 6.}
\item
Suppose~\(f\) is a nonconstant analytic function in~\(D\) such that \(|f(z)|\leq 1\). Let \(a=f(0)\). By considering the function
\[
g(z)=\frac{f(z)-a}{1-\overline{a}f(z)},
\]
 prove that
\[
\frac{|a|-|z|}{1+|a||z|}\leq |f(z)|\leq\frac{|a|+|z|}{1-|a||z|}
\]
 for \(|z|<1\).
\def\ProblemHeadingText{Problem 7.}
\item
This problem has two parts.
\begin{examsubparts}
\item[\textnormal{(a)}]
Does there exist an analytic mapping \(f:D\to D\) such that \(f(0)=0\) and \(f(i/4)=1/3\)? Justify your answer.
\item[\textnormal{(b)}]
Does there exist an analytic mapping~\(f\) of \(\{z\in\mathbb{C}\mid \operatorname{Re}(z)>0\}\) into itself such that \(f(3)=3\) and \(f(9)=6\)? Justify your answer.
\end{examsubparts}
\def\ProblemHeadingText{Problem 8.}
\item
Let~\(\mathcal{F}\) be the collection of all analytic mappings of~\(D\) into \(\{z\in\mathbb{C}\mid \operatorname{Re}(z)>0\}\) such that \(f(0)=1\). Show that~\(\mathcal{F}\) is a normal family.
\end{examproblems}
\end{document}
